% 第3章 3.3.3：积分核与伴随核关于对角线互为镜像。
% (a) (Vf)(x)=int_0^1 K(x,t)f(t)dt，K(x,t)=1_{t<x}：纵轴为输出变量 x，横轴为输入变量 t，阴影为 K=1。
% (b) 伴随 (V*g)(t)=int_0^1 K*(t,x)g(x)dx，K*(t,x)=conj K(x,t)=1_{x>t}：纵轴为输出变量 t，横轴为输入变量 x。
\begin{tikzpicture}[font=\figurefont\footnotesize,scale=2.6]
\begin{scope}
  \fill[FigureBlueFill] (0,0) -- (0,1) -- (1,1) -- cycle;
  \draw[FigureStroke,line width=.6pt] (0,0) rectangle (1,1);
  \draw[FigureMuted,dashed,line width=.5pt] (0,0) -- (1,1);
  \node[FigureBlue] at (0.3,0.72) {$K=1$};
  \node[FigureMuted] at (0.72,0.3) {$K=0$};
  \node[below] at (0,0) {$0$}; \node[below] at (1,0) {$1$}; \node[left] at (0,1) {$1$};
  \node[below] at (0.5,-0.04) {输入$t$};
  \node[left,rotate=90,anchor=south] at (-0.07,0.5) {输出$x$};
  \node[anchor=south] at (0.5,1.04) {(a) $K(x,t)=\mathbf 1_{\{t<x\}}$};
\end{scope}
\begin{scope}[xshift=1.5cm]
  \fill[FigureGreenFill] (0,0) -- (1,0) -- (1,1) -- cycle;
  \draw[FigureStroke,line width=.6pt] (0,0) rectangle (1,1);
  \draw[FigureMuted,dashed,line width=.5pt] (0,0) -- (1,1);
  \node[FigureGreen] at (0.72,0.3) {$K^*=1$};
  \node[FigureMuted] at (0.3,0.72) {$K^*=0$};
  \node[below] at (0,0) {$0$}; \node[below] at (1,0) {$1$}; \node[left] at (0,1) {$1$};
  \node[below] at (0.5,-0.04) {输入$x$};
  \node[left,rotate=90,anchor=south] at (-0.07,0.5) {输出$t$};
  \node[anchor=south] at (0.5,1.04) {(b) $K^*(t,x)=\overline{K(x,t)}$};
\end{scope}
\end{tikzpicture}
